En este laboratorio se analizó el rendimiento práctico de algoritmos de ordenamiento (MergeSort, QuickSort, PatienceSort, std::sort) y de multiplicación de matrices cuadradas (Naive, Strassen). Mediante implementaciones en C++ y scripts de automatización, se midieron los tiempos de ejecución para distintos tamaños y distribuciones de datos. El objetivo principal fue contrastar la complejidad asintótica teórica de cada algoritmo con su comportamiento real en la máquina. Los resultados demostraron que los detalles de implementación impactan drásticamente el rendimiento, destacando la superioridad de \texttt{std::sort} para arreglos grandes y del algoritmo Naive por sobre Strassen en matrices estándar.
